\documentclass[preprint, 3p,
authoryear]{elsarticle} %review=doublespace preprint=single 5p=2 column
%%% Begin My package additions %%%%%%%%%%%%%%%%%%%

\usepackage[hyphens]{url}

  \journal{Journal of Cycling and Micromobility
Research?} % Sets Journal name

\usepackage{graphicx}
%%%%%%%%%%%%%%%% end my additions to header

\usepackage[T1]{fontenc}
\usepackage{lmodern}
\usepackage{amssymb,amsmath}
% TODO: Currently lineno needs to be loaded after amsmath because of conflict
% https://github.com/latex-lineno/lineno/issues/5
\usepackage{lineno} % add
  \linenumbers % turns line numbering on
\usepackage{ifxetex,ifluatex}
\usepackage{fixltx2e} % provides \textsubscript
% use upquote if available, for straight quotes in verbatim environments
\IfFileExists{upquote.sty}{\usepackage{upquote}}{}
\ifnum 0\ifxetex 1\fi\ifluatex 1\fi=0 % if pdftex
  \usepackage[utf8]{inputenc}
\else % if luatex or xelatex
  \usepackage{fontspec}
  \ifxetex
    \usepackage{xltxtra,xunicode}
  \fi
  \defaultfontfeatures{Mapping=tex-text,Scale=MatchLowercase}
  \newcommand{\euro}{€}
\fi
% use microtype if available
\IfFileExists{microtype.sty}{\usepackage{microtype}}{}
\usepackage[]{natbib}
\bibliographystyle{elsarticle-harv}

\ifxetex
  \usepackage[setpagesize=false, % page size defined by xetex
              unicode=false, % unicode breaks when used with xetex
              xetex]{hyperref}
\else
  \usepackage[unicode=true]{hyperref}
\fi
\hypersetup{breaklinks=true,
            bookmarks=true,
            pdfauthor={},
            pdftitle={Exploring attitudes towards shared electric mobility hubs (eHUBs) through stakeholder and community workshops},
            colorlinks=false,
            urlcolor=blue,
            linkcolor=magenta,
            pdfborder={0 0 0}}

\setcounter{secnumdepth}{5}
% Pandoc toggle for numbering sections (defaults to be off)


% tightlist command for lists without linebreak
\providecommand{\tightlist}{%
  \setlength{\itemsep}{0pt}\setlength{\parskip}{0pt}}

% From pandoc table feature
\usepackage{longtable,booktabs,array}
\usepackage{calc} % for calculating minipage widths
% Correct order of tables after \paragraph or \subparagraph
\usepackage{etoolbox}
\makeatletter
\patchcmd\longtable{\par}{\if@noskipsec\mbox{}\fi\par}{}{}
\makeatother
% Allow footnotes in longtable head/foot
\IfFileExists{footnotehyper.sty}{\usepackage{footnotehyper}}{\usepackage{footnote}}
\makesavenoteenv{longtable}



\usepackage{booktabs}
\usepackage{longtable}
\usepackage{array}
\usepackage{multirow}
\usepackage{wrapfig}
\usepackage{float}
\usepackage{colortbl}
\usepackage{pdflscape}
\usepackage{tabu}
\usepackage{threeparttable}
\usepackage{threeparttablex}
\usepackage[normalem]{ulem}
\usepackage{makecell}
\usepackage{xcolor}



\begin{document}


\begin{frontmatter}

  \title{Exploring attitudes towards shared electric mobility hubs
(eHUBs) through stakeholder and community workshops}
    \author[]{James Reynolds%
  \corref{cor1}%
  \fnref{1}}
   \ead{james.reynolds@atu.ie} 
    \author[]{Noreen Armstrong%
  %
  \fnref{2}}
   \ead{noreen.armstrong@atu.ie} 
    \author[]{Amaya Vega%
  %
  \fnref{2}}
   \ead{amaya.vega@atu.ie} 
    \author[]{Brian McCann%
  %
  \fnref{1}}
   \ead{brian.mccann@atu.ie} 
    \author[]{Sarbast Molsem%
  %
  \fnref{3}}
   \ead{molsem@tcd.ie} 
    \author[]{Agnieszka Stefaniec%
  %
  \fnref{4}}
   \ead{a.stefaniec@soton.ac.uk} 
    \author[]{Brian Caulfield%
  %
  \fnref{3}}
   \ead{brian.caulfield@tcd.ie} 
      \cortext[cor1]{Corresponding author}
    \fntext[1]{Atlantic Technological University, Sligo campus, Ireland}
    \fntext[2]{Atlantic Technological University, Galway City campus,
Ireland}
    \fntext[3]{Centre for Transport Research, Dept of Civil, Structural
and Environmental Engineering, Trinity College Dublin, Ireland}
    \fntext[4]{Department of Decision Analytics and Risk at Southampton
Business School, University of Southampton, UK}
  
  \begin{abstract}
  This is the abstract.

  It consists of two paragraphs.
  \end{abstract}
    \begin{keyword}
    mobility hub \sep eHUBs \sep 
    co-design
  \end{keyword}
  
 \end{frontmatter}

\section{Introduction}\label{introduction}

The release of Toronto's \emph{Big Move} in 2008 \citep{Salsberg:2010aa}
provides an early example of the `mobility hub' concept, in the context
of regional transport and landuse planning. More recently, the ``Shared
Mobility Hub'' idea has been introduced, wherein ``\ldots emphasis is
given on shared mobility options such as bikes, scooters and cars(,
while c)onnection with public transport (transit) is desirable but not
necessary'' \citep{Roukouni:2023aa}. Shared \emph{electric} mobility
hubs, sometimes called eHUBs \citep{Bosehans:2023aa}, add the
possibility of on-site recharging of micro-mobility vehicles (as well as
shared e-cars). Along such lines, four shared e-mobility hubs have
recently been implemented in Ireland as part of the Electric shaRed
mOBilityhUbS Trial (ROBUST) project \citep{ROBUST:2025aa}. Each hub
includes e-bikes, e-cargo bikes, and shared electric cars as part of the
`base' offering, as well as power infrastructure to support recharging
of these vehicles, and other potential future uses
\citep[\citet{Atlantic-Technological-University:2025aa},
\citet{Waterford-City-and-County-Council:2025aa}]{Atlantic-Technological-University:2025ab}.

Shared micromobility, including dockless e-bike and e-scooter systems,
have variously faced challenges, including: poor user behaviour with
respect to parking; a need for workers to re-position vehicles and
replace batteries; poor build quality of many such vehicles; market
distortion, including related to large fleets in central areas; and the
risk of fires when battery recharging is done using non-bespoke or
appropriately designed facilities (e.g.~in gig-workers' houses)(see
\citet{Zhang:2022aa}, \citet{Meigs:2023aa}, \citet{Ferguson:2025aa} and
others). In contrast, at the ROBUST hubs there is a dedicated charging
and parking space for each micro-mobility vehicle, which has to be
returned to the same hub from which it was hired (i.e.~A-to-A trips
only). This may provide operational advantages (minimal staffing), but
also for a more community-based service profile (with hubs potentially
located in small towns or outside of city centres) so as to help towards
providing for last-mile and and more occasional trips that cannot be
easily accomplished on foot, personal bicycle or by transit. A nearby
eHUB might also help communities transition to electric cars, without
the large capital expense involved in purchasing a private e-car.

As of yet, however, not much is known about eHUBS with on-site
recharging of micro-mobility vehicles. It is unclear how the
return-to-origin (A-to-A) restriction might impact on users and
useability/practicality.\\
Nor is it clear what potential hub offerings might better meet community
needs (including through expansion). There also appears to be only
limited understandings about how best to communicate information about
such a hub and how to use it.

Research understandings about eHUBs appear to be at an early stage, for
which many authors\footnote{See, for example, \citet{Bonoma:1984aa},
  \citet{Benbasat:1987aa}, \citet{Eisenhardt:1989aa},
  \citet{Meredith:1993aa}, @ Cavaye:1996aa, \citet{Darke:1998aa},
  \citet{Stuart:2002aa}, \citet{Voss:2002aa} and \citet{Yin:2018ab}.}
suggest adopting an exploratory case research approach.
Industrial-design-based approaches, including those well suited to
transport research \footnote{See \citet{Coxon:2018aa},
  \citet{Coxon:2021ab}, \citet{Gaitan:2021aa}, \citet{Napper:2021aa},
  \citet{Pink:2021aa}, \citet{UK-Design-Council:0aa} and others.}, such
variously adopting iterative, co-design and participatory, and creative
processes. As such, this paper reports the findings from workshops with
members of the public and project stakeholders held in Letterkenny,
Galway and Waterford to explore attitudes, needs and ideas as to how the
`base' shared e-mobility hub might be leveraged and enhanced so as to
better meet each community's transport and other needs. The cross-case
comparisons reported in this paper might provide policy-makers,
practitioners and researchers further insights into the types of
offerings that might be provided at shared e-mobility hubs and other
transport facilities in a variety of contexts (small town, secondary
town centre, suburban locations etc.) so as to enhance sustainable
mobility and community outcomes in future such transport facility
developments.

The rest of the paper is structured as follows: INSERT STRUCTURE HERE

\section{Background and context}\label{background-and-context}

\subsection{Shared micro-mobility
problems}\label{shared-micro-mobility-problems}

\subsection{Return-to-origin (A-to-A) restriction might impact on users
and
useability/practicality.}\label{return-to-origin-a-to-a-restriction-might-impact-on-users-and-useabilitypracticality.}

\subsection{Potential hub offerings}\label{potential-hub-offerings}

THAT EXPLORATORY MODULAR THING

\subsection{Communicate information about such a hub and how to use
it,}\label{communicate-information-about-such-a-hub-and-how-to-use-it}

SEAMLESS

\subsection{Context of the Electric shaRed mOBilityhUbS Trial (ROBUST)
project}\label{context-of-the-electric-shared-mobilityhubs-trial-robust-project}

The Electric shaRed mOBility hUbS Trial (ROBUST) project involves four
sites: near the Westside Library in Galway; in Justic

\begin{longtable}[]{@{}
  >{\raggedright\arraybackslash}p{(\linewidth - 8\tabcolsep) * \real{0.1200}}
  >{\raggedright\arraybackslash}p{(\linewidth - 8\tabcolsep) * \real{0.1520}}
  >{\raggedright\arraybackslash}p{(\linewidth - 8\tabcolsep) * \real{0.1360}}
  >{\raggedright\arraybackslash}p{(\linewidth - 8\tabcolsep) * \real{0.2560}}
  >{\raggedright\arraybackslash}p{(\linewidth - 8\tabcolsep) * \real{0.3360}}@{}}
\caption{ROBUST hub characteristics}\tabularnewline
\toprule\noalign{}
\begin{minipage}[b]{\linewidth}\raggedright
Characteristic
\end{minipage} & \begin{minipage}[b]{\linewidth}\raggedright
Letterkenny
\end{minipage} & \begin{minipage}[b]{\linewidth}\raggedright
Galway
\end{minipage} & \begin{minipage}[b]{\linewidth}\raggedright
Waterford
\end{minipage} & \begin{minipage}[b]{\linewidth}\raggedright
Dublin
\end{minipage} \\
\midrule\noalign{}
\endfirsthead
\toprule\noalign{}
\begin{minipage}[b]{\linewidth}\raggedright
Characteristic
\end{minipage} & \begin{minipage}[b]{\linewidth}\raggedright
Letterkenny
\end{minipage} & \begin{minipage}[b]{\linewidth}\raggedright
Galway
\end{minipage} & \begin{minipage}[b]{\linewidth}\raggedright
Waterford
\end{minipage} & \begin{minipage}[b]{\linewidth}\raggedright
Dublin
\end{minipage} \\
\midrule\noalign{}
\endhead
\bottomrule\noalign{}
\endlastfoot
Location & Justice Walsh Road & Westside & Ballybricken & Dundrum \\
Context & Town centre & Secondary centre & Town centre / shopping
precinct & Close to a major suburban shopping centre \\
Nearby transit & Bus station & Bus stops & Railway and bus stops & Light
Rail Transit and bus stops \\
Pop nearby* & 5,656 & 10,146 & 11,919 & 13,542 \\
Pop rest\# & 17,042 & 76,186 & 48,504 & 1,256,736 \\
Launch date & 5/9/25 & 21/3/25 & 31/10/25 & 30/10/26 \\
\end{longtable}

\subsubsection{Westside, Galway}\label{westside-galway}

\subsubsection{Ju, Galway}\label{ju-galway}

Letterkenny, Waterford and Dublin, where `initial offers' of four
electric bikes, one electric cargo bike and two electric cars were
implemented during 2025 and 2026. These `initial offers' were largely
designed and implemented by engineers, planners and other technical
specialists to meet the basic needs and functionality requirements
associated with parking and charging the seven vehicles at each site.
Design approval processes with local authorities, to meet safety
requirements and to respond to other technical necessities were
undertaken, in part relying on those developed for the implementation of
shared electric car, docked bike and similar technologies on an
individual vehicle or mode basis (i.e.~not as part of a `mobility hub').
There are, however, a wide range of research questions related to what
users might want, need, expect or otherwise desire from a `mobility
hub', and the extent to which these are fulfilled by the `initial
offers', as implemented, is unclear.

\citet{Yin:2018ab}

Overall, therefore, this exploratory study focuses on three questions
that appear especially relevant to eHUBs where all vehilces (including
micro-mobility vehicles) have their own dedicated parking space and
charging infrastructure. Taking the ROBUST hubs as initially implemented
as a `base scenario', the study examines:

\begin{itemize}
\tightlist
\item
  How a return-to-origin (A-to-A) restriction minght impact on users and
  useability?
\item
  How hub offerings might better meet community needs (including through
  expansion)?
\item
  How information about a hub and how to use it be effectively
  communicated?
\end{itemize}

\section{Methodology}\label{methodology}

\subsection{Study design}\label{study-design}

This study adopts a case-research design, based around three of the four
hubs in the Electric shaRed mOBilityhUbS Trial (ROBUST) project. Case
studies can be particularly useful for exploratory research, where the
researchers do not have control over events (limiting the feasibility of
case-control experiments), and where questions relate to `how' or `why'
(\citet{Yin:2018ab} and others). While quantitative approaches (based
around having large, random and independent samples, and allowing
statistically significant results to be obtained for research relating
to `how many' or `what' types of questions) case research in contrast
can allow for greater depth of engagement with each individual unit and
its context. Some misunderstandings about case-based methods, however,
relate to the sampling frame, in that cases might not be selected
randomly or because they are representative. Instead, the methodological
literature emphasises that case researchers should consider theoretical
reasons in selecting cases, for instance because they represent a
critical, leading or insightful case, or might provide for replication
or the examination of polar opposites. The `duality criterion' is also a
consideration, in that case researchers need to be at the same time
grounded in the context of each case, allowing phenomenon to be examined
in great depth, yet at the same time seeking generalisation such that
findings can be reasonably extrapolated as applying to other
circumstances beyond the studied cases (\citet{Benbasat:1987aa},
\citet{Bonoma:1984aa}, \citet{Darke:1998aa}, \citet{Meredith:1998aa},
\citet{Stuart:2002aa}, \citet{Denscombe:2007aa},
\citet{Ketokivi:2014aa}, \citet{Yin:2018ab} and others)

Here the selection of the three included cases is largely forced
(i.e.~for non-theoretical reasons) due to the circumstances of the
larger ROBUST project, in that the Dublin hub implementation had been
delayed for reasons outside the researchers control. Likewise, site
locations and many other f

the conditions of research funding. However, this also represents a
unique opportunity, in that the four implemented hubs exist because of
this funding and that there are no known other locations where such hubs
are being trialed or where researchers might have the same sort of
access to data, stakeholders or other aspects of the project. Research
questions here also relate to how user-centric design, co-design and
similar approaches might be used to improve outcomes, support mobility
and community at individual sites, and in understanding how and why some
user-centric design approaches might be more or less effective in the
context of mobility hubs. Tools developed here might be generalisable to
other locations, such that those implementing future hubs might have
resources that can be used for maximising the potential of shared
electric mobility at each location. As well, findings about how, why and
what potential additions or modifications to the `initial offers'
implemented at Letterkenny, Galway, Waterford and Dundrum may provide
insights into how such `initial offers' might be varied, improved or
otherwise enhanced, including through the potential development of
different `initial offers' for different contexts. Hence, findings from
research at each of the ROBUST sites might provide contrasts about how
needs, desires and the like vary between larger cities and smaller
towns, central versus suburban contexts, and likewise across the range
of characteristics covered by the four sites. These sites, their
surrounding contexts and the `initial offers' are each are described in
more detail in the following. The `Challenge' for this research (c.f. UK
Design Council's Framework for Innovation above) is effectively set by
these sites and implementations, as initially implemented, in that the
outcome objective is to: \emph{discover} the initial offer's strengths,
weaknesses and opportunities for imporvement; \emph{define} these as a
problem that can be addressed through industrial design, user-centric,
participatory or similar methods; \emph{develop} options; and
\emph{deliver} an improved `initial offer' that might form the basis for
future mobility hub implementations (including at scale).

In-depth exploration of how the ROBUST project was developed and
implemented as a piece of `transport policy' is provided in
\citet{Reynolds:0ab}, while \citet{Reynolds:0aa} details the processes
used to identify options, evaluate and select between locations
themselves. The hubs were launched in: March 2025 in Galway
\citep{Atlantic-Technological-University:2025ab}; September 2025 in
Letterkenny \citep{Atlantic-Technological-University:2025aa}; November
2025 in Waterford \citet{Waterford-City-and-County-Council:2025aa}; and
XXX 2026 in Dundrum. The initial offer at each hub includes two
components: an electric bicycle / cargo-bicycle parking and charging
faciliy; and two electric car parking spaces with a charging facility.

While members of the ROBUST research team had been involved in the
development, siting and other aspects of implementation prior to the
launch of each hub, by happenstance there was a switch-over of personnel
through 2025 and 2026, with the lead author starting in September 2025,
which provides a useful demarcation point for the purposes of this case
study. Hence, September 2025 is defined as the commencement of the
`post-initial-implementation co-design-related activities' as it
represents a transition (and handover) from many of those involved in
the development of the `initial offer' and to this research. The
following table maps research activities, in chronological order, to
each of the Double Diamond design phases.

\subsection{Participants}\label{participants}

\begin{longtable}[t]{lrllr}
\caption{\label{tab:workshop_participants}Summary of workshops}\\
\toprule
Workshop & Date & Location & Type & Participants\\
\midrule
1 & 23/09/26 & Galway & Public* & 11?\\
2 & 24/09/26 & Galway & Public & 0\\
3 & 07/10/26 & Waterford & Public & 2\\
4 & 08/10/26 & Waterford & Stakeholders & 3\\
5 & 13/10/26 & Letterkenny & Stakeholders & 4?\\
\addlinespace
6 & 29/10/26 & Galway & Stakeholders & 5\\
\bottomrule
\end{longtable}

\subsection{Material and Instruments}\label{material-and-instruments}

\subsection{Procedure}\label{procedure}

\subsection{Analysis}\label{analysis}

\section{Results}\label{results}

\section{Discussion}\label{discussion}

\subsection{Limitiations}\label{limitiations}

\subsubsection{Directions for furture
research}\label{directions-for-furture-research}

\section{Conclusions}\label{conclusions}

\emph{Declaration of competing interests:} TO BE COMPLETED\footnote{See
  https://www.sciencedirect.com/journal/journal-of-cycling-and-micromobility-research/publish/guide-for-authors\#writing-and-formatting-file-format}

\emph{Funding:} This work was supported by the Sustainable Energy Agency
of Ireland {[}Electric shaRed mOBility hUbS Trial (ROBUST), {]} and
involved: Trinity College Dublin; Atlantic Technological Uuniversity
(ATU); ESB; and Enterprise Holdings.

\begingroup\fontsize{8}{10}\selectfont

\begin{longtable}[t]{>{\raggedright\arraybackslash}p{3cm}>{\raggedright\arraybackslash}p{6cm}>{\raggedright\arraybackslash}p{0.5cm}>{\raggedright\arraybackslash}p{0.5cm}>{\raggedright\arraybackslash}p{0.5cm}>{\raggedright\arraybackslash}p{0.5cm}>{\raggedright\arraybackslash}p{0.5cm}>{\raggedright\arraybackslash}p{0.5cm}>{\raggedright\arraybackslash}p{0.5cm}}
\caption{\label{tab:contribution_matrix_template}Author contribution matrix, as per CRediT (Contribution Roles Taxonomy}\\
\toprule
Component & Description & NA & BC & BM & SM & JR & AS & AV\\
\midrule
Conceptualization & Ideas; formulation or evolution of overarching research goals and aims. &  &  &  &  & X &  & \\
Data Curation & Management activities to annotate (produce metadata), scrub data and maintain research data (including software code, where it is necessary for interpreting the data itself) for initial use and later reuse. &  &  &  &  &  &  & \\
Formal Analysis & Application of statistical, mathematical, computational, or other formal techniques to analyze or synthesize study data. &  &  &  &  &  &  & \\
Funding Acquisition & Acquisition of the financial support for the project leading to this publication. &  & X &  &  &  & X & \\
Investigation & Conducting a research and investigation process, specifically performing the experiments, or data/evidence collection. & X &  &  &  & X &  & \\
\addlinespace
Methodology & Development or design of methodology; creation of models. &  &  &  &  & X &  & \\
Project Administration & Management and coordination responsibility for the research activity planning and execution. & X &  &  & X &  &  & \\
Resources & Provision of study materials, reagents, materials, patients, laboratory samples, animals, instrumentation, computing resources, or other analysis tools. &  &  &  &  &  &  & \\
Software & Programming, software development; designing computer programs; implementation of the computer code and supporting algorithms; testing of existing code components. &  &  &  &  &  &  & \\
Supervision & Oversight and leadership responsibility for the research activity planning and execution, including mentorship external to the core team. &  &  & X &  &  &  & X\\
\addlinespace
Validation & Verification, whether as a part of the activity or separate, of the overall replication/reproducibility of results/experiments and other research outputs. &  &  &  &  &  &  & \\
Visualization & Preparation, creation and/or presentation of the published work, specifically visualization/data presentation. &  &  &  &  &  &  & \\
Writing – Original Draft Preparation & Creation and/or presentation of the published work, specifically writing the initial draft (including substantive translation). &  &  &  &  & X &  & \\
Writing – Review \& Editing & Preparation, creation and/or presentation of the published work by those from the original research group, specifically critical review, commentary or revision – including pre- or post-publication stages. &  &  &  &  &  &  & \\
\bottomrule
\end{longtable}
\endgroup{}

\renewcommand\refname{References}
\bibliography{01_exploratory_workshops.bib}


\end{document}
